\section{Contoh}
\label{sec:contoh-bab14}

Dua contoh berikut memperlihatkan prosedur berparameter dan macro.
Contoh pertama mengirim parameter melalui stack. Contoh kedua
memakai macro dengan parameter karakter.

\subsection{Contoh 1: Prosedur Berparameter melalui Stack}

Program pada Kode~\ref{lst:prosedur-luas} menghitung luas persegi
panjang melalui prosedur \code{HitungLuas}.

\lstinputlisting[
    language=MASM,
    caption={Program \file{prosedur-luas.asm} memakai prosedur
    berparameter melalui stack.},
    label={lst:prosedur-luas}
]{\CodeDir/prosedur-luas.asm}

Bagian penting program:

\begin{itemize}
    \item Pemanggil mendorong \code{tinggi} lebih dahulu, lalu
          \code{lebar}. Parameter yang didorong terakhir berada paling
          dekat dengan puncak stack.
    \item \asm{CALL HitungLuas} menyimpan alamat kembali ke stack.
    \item Prosedur menyimpan \reg{BP} lama dengan \asm{PUSH BP}, lalu
          menyalin \reg{SP} ke \reg{BP}.
    \item \lstasm|MOV AX, [BP+4]| membaca parameter pertama, yaitu
          \code{lebar}.
    \item \lstasm|MUL WORD PTR [BP+6]| mengalikan dengan parameter
          kedua, yaitu \code{tinggi}.
    \item \asm{RET} \code{4} kembali ke pemanggil sekaligus membuang
          empat byte parameter.
\end{itemize}

Hasilnya adalah $5 \times 10 = 50$ atau \hex{0032}, disimpan pada
variabel \code{luas} dan register \reg{AX}.

\subsection{Contoh 2: Macro dengan Parameter}

Program pada Kode~\ref{lst:macro-demo} mendefinisikan macro
\code{Cetak\_Kar} yang mencetak satu karakter sebanyak tiga kali.

\begin{lstlisting}[
    language=MASM,
    caption={Macro \code{Cetak\_Kar} dan pemakaiannya.},
    label={lst:macro-demo}
]
.MODEL SMALL
.STACK 100h

Cetak_Kar MACRO Kar
    MOV CX, 3
    MOV AH, 02h
    MOV DL, Kar
Ulang:
    INT 21h
    LOOP Ulang
ENDM

.CODE
START:
    Cetak_Kar 'S'       ; cetak 'S' tiga kali

    MOV AH, 4Ch
    MOV AL, 00h
    INT 21h
END START
\end{lstlisting}

Saat assembler menemukan \lstasm|Cetak_Kar 'S'|, isi macro disalin ke
tempat itu dengan setiap kemunculan \code{Kar} diganti \code{'S'}.
Keluaran program adalah \code{SSS}.

\begin{lstlisting}[
    style=sessiongaya,
    caption={Keluaran program macro.},
    label={lst:keluaran-macro}
]
SSS
\end{lstlisting}

\begin{infobox}
\vspace{1.1em}
Macro pada contoh ini memakai label \code{Ulang}. Label tersebut
aman karena macro hanya dipakai sekali. Bila macro dipakai lebih dari
sekali, label yang sama akan disalin berulang dan assembler dapat
melaporkan kesalahan label ganda.
\end{infobox}
